\subsection{Every first coordinate unramified outside three}

\begin{statement}{E.51}{Proved}{All nonzero three-power Kummer first coordinates in the inherited quotient}
Retain the based arithmetic quotient $H$ of E.46 and the
positive real root system defining $\kappa=\kappa_3(3)$ in
E.49. Let $r\in\mathbb Z_3\setminus\{0\}$, and write
\[
 r=3^k e,\qquad k=v_3(r),\quad e\in\mathbb Z_3^\times,
 \qquad M=3^{k+1},\quad h=(M-1)/2.
 \tag{E51.1}
\]
Define a set of rational primes and its restricted
three-adic product by
\[
 \begin{split}
 \mathscr P_k&=\{p\ne3:
       v_3(\operatorname{ord}_{\mathbb F_p^\times}(3))
                                                \le k+1\},\\
 \mathcal R(\mathscr P_k)&=
 \{(b_p)\in\prod_{p\in\mathscr P_k}\mathbb Z_3:
       \text{for every }a,\ b_p\in3^a\mathbb Z_3
                       \text{ for almost all }p\}.
 \end{split}                                                   \tag{E51.2}
\]
The order at $p=2$ is one. After the choices of
prime norm lifts specified below, the complete global
section fiber of $H$ with first class $[3r\kappa]$
has the following additive topological coordinates:
\[
 \boxed{\quad
 \mathbb Z_3\times\mathcal R(\mathscr P_k)
                 \times\mathbb Z_3^{h}\times\mathbb Z_3[[Z]].
 \quad}                                                        \tag{E51.3}
\]
Write its coordinates as $(b_3,(b_p),(z_i)_{i=1}^h,f)$.
The scaled toric map of E.46 is exactly
\[
 (b_3,(b_p),(z_i),f)\longmapsto
 \left(r\kappa,\ b_3\kappa+
                         \sum_{p\in\mathscr P_k}b_p\kappa_3(p)\right).
 \tag{E51.4}
\]
Thus every toric fiber over an allowed second coordinate
is explicitly $\mathbb Z_3^h\times\mathbb Z_3[[Z]]$.
All the displayed groups are three-torsion-free.

Here are actual cocycle coordinates for the unramified
part. Use the Frobenius basis $\tau,\tau'=c\tau c$
of E.49, for which $\chi(\tau)=7$ and
$\kappa(\tau)=1$. Put $T=1+U$,
$W=(T-T^{-1})/2$, and
\[
 \eta_{M,i}=1-\frac{T^{-i}+T^{i-M}}2
                           \qquad(1\le i\le h).
 \tag{E51.5}
\]
For the normalized first cocycle $A=M\kappa$, every
second-cocycle class unramified outside three has
exactly one representative of the form
\[
 \begin{split}
 v(c)&=0,\\
 v(\tau)&=3b_3+\sum_{i=1}^{h}z_i\eta_{M,i}
                                +T^M W f(W^2),\\
 v(\tau')&=-v(\tau)(T^{-1}).
 \end{split}                                                   \tag{E51.6}
\]
For the prescribed cocycle $3r\kappa=eM\kappa$,
replace every occurrence of $T$ in the module
coordinate by $T^e$. The exact map accomplishing
this change is the continuous arithmetic automorphism
\[
 \Phi_e:H\xrightarrow{\sim}H,
       \qquad (a,v(T))\longmapsto(ea,v(T^e)).
 \tag{E51.7}
\]
Its inverse is $\Phi_{e^{-1}}$. It preserves $J=(3,U)$,
commutes with the based Galois action, and preserves
the second scaled toric coordinate.
On every finite root set the corresponding equivariant
bijection is $i\mapsto ei$. Pullback gives an isomorphism
from the algebra for $eM\kappa$ to the one for $M\kappa$;
its inverse transports the chosen norm systems below.
These isomorphisms commute with every tower map and norm.

For each $p\in\mathscr P_k$, choose the norm system
$\pi_{k,p}$ constructed in the proof. Its divisor
is the compatible degree-one prime above $p$, and
all other valuations, including those at three,
are zero. Its level-zero norm is $p$. The ramified
factor in the norm coordinates of E.46 is
$\prod_p\pi_{k,p}^{\,3b_p}$. Together with the
Kummer--Shapiro image of (E51.6), this specifies the
map in (E51.3), including all infinite products.

The first classes unramified outside three are
exactly $[3r\kappa]$ for $r\in\mathbb Z_3$.
E.47 supplies the fiber at zero, and (E51.3) supplies
every nonzero fiber. These are section classes of
the specified target $H$; the maps from the
intrinsic groups retain the kernels of E.46--E.47.

The program \texttt{e\_unramified\_first.py} checks
the exact coordinate maps, finite affine orbits,
field factors and norm transitions used below.

\emph{Manuscript: \S51.16, folios733v--736v,
scans172--178, (371)--(379), Part E(c),(d) and
the fourth exactness follow-up.}
\end{statement}
\begin{proof}
\emph{The change of first cocycle.}
The multiplication in $H$ is
$(a,v)(b,w)=(a+b,v+T^a w)$. Substituting $T^e$
in this identity proves that (E51.7) is a group
homomorphism. Its inverse follows by substituting
$T^{e^{-1}}$. The series $T^e-1$ has unit linear
coefficient, hence the substitution is a continuous
automorphism of $\mathbb Z_3[[U]]$ carrying $J$ to
itself. It commutes with $T\mapsto T^{\chi_g}$
and the scalar factor $\chi_g$. Thus it is an
arithmetic automorphism, induces a bijection on
the stated section fibers, and changes their first
coordinate by $e$. Augmentation is unchanged,
which proves the assertion about the second toric
coordinate. It suffices to compute the fiber for
the actual cocycle $A=M\kappa$.

\emph{Integral computation of the unramified group.}
The free arithmetic basis and the full-to-pro-three
cohomology comparison are those proved in E.49.
They apply to every finite coefficient quotient
here: over $K=\mathbb Q(\zeta_3)$ its arithmetic
image is a three-group. Its filtration by trivial
$\mathbb F_3$ modules gives vanishing of unramified
$H^2$ and $H^3$; its $H^1$ is finite. These facts
also justify the inverse-limit comparisons below.

After normalizing $v(c)=0$, write $u=v(\tau)$.
The other free value is $Cu=-u(T^{-1})$.
For $\rho=\tau'\tau^{-1}$ the coefficient operators
are
\[
 S u=7T^M u(T^7),\qquad R u=T^{-2M}u,
             \qquad C u=-u(T^{-1}).
 \tag{E51.8}
\]
They satisfy $CSC=RS$ and $CRC=R^{-1}$, matching
the actual conjugation relations of $c,\tau,\rho$.
The remaining coboundaries change $u$ by
$Dm=(S-1)m$, where $Cm=m$. Such $m$ has
augmentation zero and belongs to $J$.

Let $n=T^{-2M}-1$. The exact factorization
\[
 n=-T^{-2M}(T^M+1)(T^M-1)
 \tag{E51.9}
\]
gives $(n)=(T^M-1)$, since $T^M+1$ has constant
term two. The quotient is the finite free ring
$R_M=\mathbb Z_3[C_M]$. Write $\bar u_j$ for
the coefficient of $T^j$ in the image of $u$,
where $0\le j<M$. The involution $T\mapsto T^{-1}$
has a fixed submodule of rank $h+1$, with basis
given by $1$ and the images of $\eta_{M,i}$.
Its projection $(1+\iota)/2$ annihilates every
$\overline{Dm}$: substitution by $T^7$ commutes
with $\iota$, and $\bar m$ is anti-invariant.

We evaluate the whole quotient by these
coboundaries. Set
\[
 \begin{split}
 a&=\operatorname{aug}u,\qquad
 z_i=-(\bar u_i+\bar u_{M-i}),\\
 u_*&=u-a-\sum_{i=1}^h z_i\eta_{M,i},\\
 m&=\frac{Cu_*-Ru_*}{n},\qquad
                  v=u_*-Dm.
 \end{split}                                                   \tag{E51.10}
\]
The image of $u_*$ in $R_M$ is anti-invariant
under inversion. Thus the numerator defining $m$
vanishes in $R_M$ and the division is integral.
The domain $\Lambda=\mathbb Z_3[[U]]$ has no
$n$-torsion. Applying $C$ to the numerator gives
$-T^{2M}$ times that numerator, while
$C(nm)=-T^{2M}nCm$. Therefore $Cm=m$.
The operator identities in (E51.8) now give
\[
 C(Dm)-R(Dm)=nm,\qquad Cv=Rv.
 \tag{E51.11}
\]
Consequently $T^{-M}v$ is anti-invariant under
inversion. The integral change of variable
$W=(T-T^{-1})/2$ has inverse
$T=W+\sqrt{1+W^2}$, with constant term one.
It identifies $\Lambda$ with $\mathbb Z_3[[W]]$
and inversion with $W\mapsto-W$. This proves
$v=T^M Wf(W^2)$ for a unique $f\in\mathbb Z_3[[Z]]$.

The coordinates are unique modulo coboundaries.
The fixed projection in $R_M$ detects every $a,z_i$.
For two representatives with the same finite
coordinates, their difference has zero
$\rho$-value. A coboundary with this property
satisfies $nm=0$, hence $m=0$. Their $f$
coordinates therefore coincide as well.
All divisions above are continuous on their
closed principal ideals: multiplication by a
nonzero series is a continuous bijection from
the compact ring onto its image. This proves
the topological normal form. Restricting from
$\Lambda$ to $J$ imposes exactly
$a\in3\mathbb Z_3$, yielding (E51.6).

\emph{The exact allowable residues.}
At $p\ne3$, the twisted module is unramified.
After removing its Tate twist, an actual
arithmetic Frobenius acts on the positions by
$x\mapsto M\kappa(F_p)+px$. Its possible fixed
point is unique. An integral fixed point is,
by the root labels in E.46, exactly a compatible
system of $3^n$th roots of $3^M$ in $\mathbb Q_p$.
Cubing is an automorphism on the pro-$p$
principal units; in the finite residue-unit
group the compatible divisibility criterion
is
\[
 3\nmid\operatorname{ord}_{\mathbb F_p^\times}(3^M)
 \quad\Longleftrightarrow\quad
 v_3(\operatorname{ord}_{\mathbb F_p^\times}(3))\le k+1.
 \tag{E51.12}
\]
This proves the prime set in (E51.2).

The integral fixed-measure argument of E.49
applies verbatim to this affine map: translate
an integral fixed point to zero, and use
multiplication by $p^2$ on every other annulus.
Each of its nonzero clopen orbits is a free
$\mathbb Z_3$ orbit. A compatible invariant
integral measure there has mass divisible
by arbitrarily high powers of three, so is
zero. With a nonintegral center, the valuation
of the displacement is $t+v_3(n)$ on every
orbit, with $t<v_3(p-1)$, and the same argument
applies on its finitely many residue balls.
Intersecting with $J$ therefore gives
\[
 (J_A(-1))^{F_p}=
 \begin{cases}
 3\mathbb Z_3\delta_{x_p},&p\in\mathscr P_k,\\
 0,&p\notin\mathscr P_k,
 \end{cases}
 \tag{E51.13}
\]
where $x_p$ denotes the unique fixed point.

The finite-coefficient localization sequence
used in E.49 is exact with zero right obstruction
because of the unramified $H^2$ vanishing already
checked. The finite unramified $H^1$ groups
satisfy Mittag--Leffler. The continuous
inverse-limit comparison therefore gives
the exact sequence
\[
 0\longrightarrow H^1(G_{\mathbb Q,\{3\}},J_A)
 \longrightarrow H^1(G_{\mathbb Q},J_A)
 \longrightarrow 3\mathcal R(\mathscr P_k)
 \longrightarrow0 .                                      \tag{E51.14}
\]
Here the residue at $p$ is its coefficient
of $\delta_{x_p}$. Positional pushforward
is ideal norm, so its total mass at level
zero forces precisely the restrictedness
in (E51.2). Conversely that restrictedness
is sufficient at every finite coefficient
and positional level. This is the actual
inverse limit of the direct sums, without
replacing an image by its closure.

\emph{Explicit norm lifts of the residues.}
We give a field construction of the splitting
of (E51.14). A cofinal range is $n=d+k+1$
with $d\ge0$. The finite etale algebra of
E.46 factors as
\[
 E_{A,n}=\prod_{j=0}^{k+1} L_{d,j},\qquad
 L_{d,0}=\mathbb Q(3^{1/3^d}),\quad
 L_{d,j}=\mathbb Q(3^{1/3^d}\zeta_{3^{d+j}})
                                  \quad(j\ge1).
 \tag{E51.15}
\]
The degrees are $3^d$ and $3^d\varphi(3^j)$.
The first polynomial is Eisenstein at three.
For the others, the element $\rho=\tau'\tau^{-1}$
in the actual unramified arithmetic group
has $\chi(\rho)=1$ and $\kappa(\rho)=-2$.
Its powers fix every three-power root of unity and
permute all $3^d$ roots above a prescribed
$3\zeta_{3^j}$. These are actual arithmetic
actions on the root torsor: its Galois
action factors through that group, since
the splitting fields are unramified outside
three. Moreover the generator in (E51.15)
contains $\zeta_{3^j}$. This proves the
asserted degree and irreducibility. Equivalently,
the polynomial identity is
\[
 X^{3^{d+k+1}}-3^{3^{k+1}}
  =(X^{3^d}-3)
       \prod_{j=1}^{k+1}
       3^{\varphi(3^j)}
               \Phi_{3^j}(X^{3^d}/3).
 \tag{E51.16}
\]
Cubing the chosen generator at $d+1$ gives
the one at $d$, so these are the actual
field maps in the norm tower.

For $j\ge1$ each vertical layer is cyclic
cubic and is unramified away from three.
At $d=0$ it starts at $\mathbb Q(\zeta_{3^j})$,
whose three-primary class group is zero
by E.47. There is a unique prime above three.
A nontrivial cyclic cubic layer unramified
also at this prime would be everywhere
unramified, contradicting that class-group
calculation. Thus this prime is totally
ramified, remains unique, and has residue
degree one. The ambiguous class-number
formula proves inductively that the next
class group again has order prime to three:
its fixed three-primary subgroup is zero,
while a nonzero finite three-group acted on
by a group of order three has nonzero fixed
points. The strongly ambiguous formula
\cite[Theorem1, formula(2)]{E_LemmermeyerAmbiguous}
also shows that the ordinary unit norm is
onto, because its index is a power of three
dividing the base class number.

For the $j=0$ factor, first adjoin
$K=\mathbb Q(\zeta_3)$. The resulting vertical
layers are cyclic cubic; the valuation
$2/3^d$ of the radical proves their degree
and their unique totally ramified prime.
The same argument starts at $K$ and proves
both claims there. Extension and norm over
the quadratic map remove the three-primary
class group downstairs, and the integral
projector $(1+c)/2$ proves surjectivity of
the completed ordinary unit norms in the
pure radical tower. All factors of (E51.15)
therefore have surjective completed divisor
maps and surjective completed unit norms.

For $p\in\mathscr P_k$, its fixed position
gives a unique compatible degree-one prime
chain in these algebras. At each level,
take a generator of a prime-to-three power
of this ideal and raise it to the inverse
of that power in the three-completion.
This gives precisely the desired divisor.
Correct each succeeding norm discrepancy
by an ordinary unit using the surjective
unit norm just proved. Induction constructs
$\pi_{k,p}$ with no other valuation,
including at three. Its norm to level zero
is exactly $p$, since the three-completion
of the ordinary rational units is zero.
This supplies the claimed explicit choice
of residue lift.

Every global norm system satisfying the
cube condition of E.46 has unique residue
coefficients $3b_p$. Divide by
$\prod_p\pi_{k,p}^{3b_p}$; the remaining
class is unramified outside three and
therefore has the unique coordinates in
(E51.6). The product converges in every
power-class and positional quotient
because $b_p$ tends to zero three-adically.
The construction and its inverse are
continuous. This proves (E51.3), its group
law and topology, and (E51.4). Changing a
chosen $\pi_{k,p}$ multiplies it by an
ordinary-unit norm system and changes
the compact coordinates by the
corresponding convergent additive sum.

Finally, the completed rational Kummer
group has the restricted prime-valuation
coordinates of E.16. Being unramified
outside three sets every coordinate
except the one at three to zero. Thus
its first-coordinate classes are exactly
the stated $3r\kappa$. The zero fiber
and the nonzero fibers are all accounted
for by their explicit maps. No lift
through the intrinsic geometric kernels
is asserted by this quotient calculation.
\end{proof}
